% TEMPLATE-MILC-v3.tex - plantilla maestra para todas las guias MILC v3
% Ver scripts/generadores/build-guia-milc.py para la lista completa de
% claves esperadas. El script valida que no queden placeholders sin
% reemplazar antes de escribir el archivo final.

\documentclass[11pt,letterpaper]{article}
\usepackage[margin=1.75cm]{geometry}
\usepackage[table]{xcolor}
\usepackage{fontspec}
\usepackage[spanish]{babel}
\usepackage{tikz}
\usepackage[most]{tcolorbox}
\usepackage{tabularx}
\usepackage{array}
\usepackage{fancyhdr}
\usepackage{titlesec}
\usepackage{ragged2e}
\usepackage{enumitem}
\usepackage{microtype}

\setmainfont{Helvetica Neue}
\setsansfont{Helvetica Neue}
\setlength{\parindent}{0pt}
\setlength{\parskip}{6pt}
\hyphenpenalty=200
\exhyphenpenalty=200
\pretolerance=400
\tolerance=2000
\setlength{\emergencystretch}{1em}
\renewcommand{\arraystretch}{1.25}
\setlist{leftmargin=5mm,itemsep=1mm,topsep=1mm}
\newcolumntype{Y}{>{\RaggedRight\arraybackslash}X}

\definecolor{milcMagenta}{HTML}{E6007E}
\definecolor{milcVino}{HTML}{5A0038}
\definecolor{milcTurquesa}{HTML}{0093A5}
\definecolor{milcVerde}{HTML}{58A923}
\definecolor{milcMostaza}{HTML}{E5B400}
\definecolor{milcOcre}{HTML}{B86600}
\definecolor{milcNegro}{HTML}{241816}
\definecolor{milcGris}{HTML}{F7F7F4}
\definecolor{milcLinea}{HTML}{E8E2DD}
\definecolor{milcGradePrimary}{HTML}{0066FF}
\definecolor{milcGradeDark}{HTML}{003D99}
\definecolor{milcGradeSoft}{HTML}{D6E8FF}
\definecolor{milcGradeAccent}{HTML}{CFE7FF}
\definecolor{milcGradeText}{HTML}{FFFFFF}
\color{milcNegro}

\pagestyle{fancy}
\fancyhf{}
\lhead{\small Tecnología e Informática · Grado 11}
\rhead{\small Guía 5 · MILC v3}
\cfoot{\small\thepage}
\renewcommand{\headrulewidth}{0pt}

\newtcolorbox{titlebox}[1]{
  enhanced,colback=#1,colframe=#1,arc=7mm,boxrule=0pt,
  left=7mm,right=7mm,top=3mm,bottom=3mm,
  before skip=8pt,after skip=8pt,
  fontupper=\sffamily\bfseries\Large\color{white}
}

\newtcolorbox{softbox}[3]{
  enhanced,colback=#2,colframe=#1,borderline west={4pt}{0pt}{#1},
  arc=2mm,boxrule=.4pt,left=4mm,right=4mm,top=3mm,bottom=3mm,
  before skip=6pt,after skip=8pt,title={#3},coltitle=white,
  colbacktitle=#1,fonttitle=\sffamily\bfseries,fontupper=\RaggedRight,
  attach boxed title to top left={xshift=3mm,yshift=-2mm},
  boxed title style={arc=2mm, boxrule=0pt}
}

\newtcolorbox{drawbox}[2]{
  enhanced,colback=white,colframe=#1,arc=4mm,boxrule=1pt,
  left=5mm,right=5mm,top=4mm,bottom=4mm,before skip=8pt,after skip=8pt,
  title={#2},coltitle=white,colbacktitle=#1,fonttitle=\sffamily\bfseries,
  fontupper=\RaggedRight,
  attach boxed title to top left={xshift=4mm,yshift=-2mm},
  boxed title style={arc=3mm, boxrule=0pt}
}

\newtcolorbox{infoband}[1]{
  enhanced,colback=#1!8,colframe=#1!50,arc=2mm,boxrule=.5pt,
  left=4mm,right=4mm,top=2mm,bottom=2mm,before skip=4pt,after skip=4pt,
  fontupper=\small\RaggedRight
}

% Puente narrativo entre secciones (contrato editorial Conéctate).
% Da continuidad: "Acabas de X. Ahora vas a Y porque Z."
% Visualmente: banda izquierda en color del grado, texto en itálica pequeña.
\newtcolorbox{puentebox}{
  enhanced,colback=milcGradeSoft,colframe=milcGradeDark,
  borderline west={3pt}{0pt}{milcGradeDark},
  arc=0mm,boxrule=0pt,left=5mm,right=4mm,top=2.5mm,bottom=2.5mm,
  before skip=4pt,after skip=8pt,
  fontupper=\itshape\small\color{milcNegro}\RaggedRight
}
\newcommand{\puente}[1]{%
  \begin{puentebox}%
    \textcolor{milcGradeDark}{\textbf{\textrightarrow}}\hspace{2mm}#1%
  \end{puentebox}%
}

\newcommand{\linead}[1]{\noindent\color{milcLinea}\rule{#1}{0.5pt}\color{milcNegro}}
\newcommand{\respuesta}[1]{\par\vspace{2mm}\foreach \n in {1,...,#1}{\noindent\color{milcLinea}\rule{\linewidth}{0.5pt}\par\vspace{5mm}}\color{milcNegro}}
\newcommand{\checkbox}{\textcolor{milcGradeDark}{\fbox{\phantom{x}}}\hspace{5pt}}

\newcommand{\phasecard}[4]{
\begin{tcolorbox}[enhanced,colback=#3,colframe=#2,arc=3mm,boxrule=.5pt,height=3.05cm,valign=center,halign=center,left=1.5mm,right=1.5mm,top=2mm,bottom=2mm]
{\sffamily\bfseries\fontsize{22}{24}\selectfont\textcolor{#2}{#1}}\par
{\sffamily\bfseries\small #4}
\end{tcolorbox}
}

\begin{document}
\thispagestyle{empty}
\begin{tikzpicture}[remember picture,overlay]
  \fill[white] (current page.south west) rectangle (current page.north east);
  \path[fill=milcGradePrimary,rounded corners=16mm]
    ([xshift=.55cm,yshift=-.55cm]current page.north west) rectangle
    ([xshift=-.55cm,yshift=3.65cm]current page.south east);

  \node[anchor=north west,text=milcGradeText] at ([xshift=2.0cm,yshift=-1.85cm]current page.north west)
    {\sffamily\bfseries\fontsize{14}{16}\selectfont GRADO};
  \node[anchor=north west,text=milcGradeText,opacity=.82] at ([xshift=2.0cm,yshift=-2.75cm]current page.north west)
    {\sffamily\bfseries\fontsize{9}{11}\selectfont SERIE GUÍAS · TECNOLOGÍA E INFORMÁTICA};

  \begin{scope}[shift={([xshift=-3.05cm,yshift=-2.55cm]current page.north east)}]
    \fill[white,opacity=.80] (-.62,-.52) rectangle (.62,.52);
    \draw[milcGradeText,line width=1pt] (-.62,-.52) rectangle (.62,.52);
    \fill[milcGradeDark] (-.42,-.38) rectangle (-.22,.06);
    \fill[milcGradeAccent] (-.10,-.38) rectangle (.10,.24);
    \fill[milcGradePrimary!60!black] (.22,-.38) rectangle (.42,.38);
  \end{scope}

  \node[anchor=west,text=milcGradeText] at ([xshift=1.9cm,yshift=-12.85cm]current page.north west)
    {\sffamily\bfseries\fontsize{124}{124}\selectfont 11};
  \draw[milcGradeText,line width=1.7pt]
    ([xshift=4.52cm,yshift=-11.68cm]current page.north west) circle (.13cm);

  \node[anchor=west,text=milcGradeText,text width=8.7cm,align=left] at ([xshift=10.35cm,yshift=-11.6cm]current page.north west)
    {
     {\sffamily\bfseries\fontsize{19}{22}\selectfont Copywriting con IA ---\\escribir con criterio,\\no a ciegas}\\[4mm]
     {\sffamily\bfseries\fontsize{23}{26}\selectfont Guía 5-11-TIC}\\[2mm]
     {\sffamily\fontsize{13}{16}\selectfont Período 1 · Presencia y marca digital}\\[5mm]
     {\sffamily\bfseries\fontsize{11}{13}\selectfont Institución Educativa Sor María Juliana}\\[1mm]
     {\sffamily\fontsize{10}{12}\selectfont Autor: Dr. Álvaro Cárdenas Orozco}
    };

  \path[fill=milcGradeSoft,rounded corners=7mm]
    ([xshift=1.35cm,yshift=.85cm]current page.south west) rectangle
    ([xshift=-1.35cm,yshift=4.25cm]current page.south east);
  \draw[milcGradeDark,line width=.65pt,rounded corners=7mm]
    ([xshift=1.35cm,yshift=.85cm]current page.south west) rectangle
    ([xshift=-1.35cm,yshift=4.25cm]current page.south east);
  \node[anchor=center,text=milcNegro,text width=17.0cm,align=center] at ([yshift=2.55cm]current page.south)
    {\sffamily\fontsize{10}{12}\selectfont Metodología MILC · Apertura ancestral · Pensamiento computacional · Triángulo Dussel-Estoicismo-Floridi\\
    \textcolor{milcGradeDark}{\bfseries Producto:} Página ``Sobre mí'' del sitio, coescrita con IA y revisada con criterio personal};
\end{tikzpicture}
\mbox{}
\newpage

\section*{Apertura: Copywriting con IA --- escribir con criterio, no a ciegas}

\begin{softbox}{milcGradeDark}{milcGradeSoft}{Saber ancestral}
La \textbf{tradición oral} de los pueblos del Pacífico, los Andes y el Valle no era una transmisión pasiva: cada generación tomaba la historia de la anterior, la adaptaba a su contexto y se la pasaba a la siguiente. Lo que llegaba a tus abuelos no era exactamente lo que sus abuelos contaron --- pero conservaba la esencia. La autoría era \emph{compartida} entre generaciones, con la voz viva del que cuenta hoy.
\end{softbox}

\begin{softbox}{milcTurquesa}{milcGris}{Saber contemporáneo}
Hoy tienes un nuevo ``ancestro colaborador'': la IA generativa (ChatGPT, Claude, Gemini). Te puede ayudar a escribir más rápido y mejor estructurado. Pero como cualquier coautor, necesita que tú aportes lo único que ella no tiene: \textbf{quién eres y por qué importa}. Sin esa voz tuya encima, la IA escribe textos pasables que no son de nadie.
\end{softbox}

\begin{softbox}{milcMostaza}{milcGris}{Pregunta-puente}
\textit{¿Cómo se hereda la voz de una generación a otra sin que se diluya su identidad? ¿Y cómo escribes con una herramienta que puede sonar a cualquiera?}
\end{softbox}

\begin{softbox}{milcVerde}{milcGris}{Saber y saber hacer}
Al terminar podrás: (1) \textbf{identificar} patrones únicos de tu voz al escribir; (2) \textbf{analizar} cómo cambia un texto cuando lo escribe la IA en lugar de ti; (3) \textbf{crear} la página ``Sobre mí'' del sitio en colaboración con IA y revisión tuya; (4) \textbf{evaluar} si el resultado final suena a ti.
\end{softbox}

\small
\begin{tabularx}{\textwidth}{>{\bfseries}p{3.0cm}Y}
\rowcolor{milcGradePrimary!12}Autor & Dr. Álvaro Cárdenas Orozco\\
DBA & Comunicación profesional en entornos digitales (MEN, Lineamientos T\&I)\\
Referentes & Floridi (infoética) · Dussel (estética de la liberación) · Estoicismo\\
Producto final & Página ``Sobre mí'' del sitio, coescrita con IA y revisada con criterio personal\\
\end{tabularx}
\normalsize

\puente{Antes de empezar, mira el viaje completo. Hoy le pones \emph{voz} a tu sitio --- escribirás la página ``Sobre mí'' en colaboración con IA, sin perder lo que solo tú puedes decir.}

\begin{titlebox}{milcGradeDark}
Ruta MILC: así vamos a aprender
\end{titlebox}
\begin{tabularx}{\textwidth}{>{\centering\arraybackslash}X>{\centering\arraybackslash}X>{\centering\arraybackslash}X>{\centering\arraybackslash}X}
\phasecard{1}{milcVerde}{milcGris}{Escuta\\[-1mm]\small Observo, escucho y pregunto.} &
\phasecard{2}{milcTurquesa}{milcGris}{Sistematización\\[-1mm]\small Pensamiento computacional.} &
\phasecard{3}{milcMagenta}{milcGris}{Praxis\\[-1mm]\small Construyo y aplico.} &
\phasecard{4}{milcVino}{milcGris}{Evaluación\\[-1mm]\small Reflexión liberadora.}
\end{tabularx}


\puente{Antes de pedirle a la IA que escriba contigo, tienes que conocer tu propia voz. ¿Qué palabras repites? ¿Qué giros usas? Solo así notarás cuándo la IA escribe como tú y cuándo escribe como nadie.}

\begin{titlebox}{milcVerde}
Fase 1 · Escuta
\end{titlebox}

\begin{softbox}{milcVerde}{milcGris}{Escena para escuchar}
\textbf{Actividad 1 · IDENTIFICA --- Tu voz al escribir} (15 min · individual). Mira los últimos 20 mensajes de WhatsApp que escribiste (no copies-pegues los que te enviaron). Identifica: ¿qué palabras o frases repites? ¿usas comas o frases cortas? ¿usas emojis, cuáles? ¿hablas con tono formal, informal, irónico? Tu voz al escribir es ese patrón.
\end{softbox}

\checkbox Revisé mis últimos 20 mensajes de WhatsApp con ojo de ``observador''.\\
\checkbox Identifiqué al menos 3 patrones repetidos en mi forma de escribir.\\
\checkbox Anoté qué tono predomina en mi voz: formal, informal, irónico, cálido.

\begin{infoband}{milcVerde}
\textbf{Lo que va al cuaderno (Actividad 1 · IDENTIFICA).} Título: «Actividad 1 --- Mi voz al escribir». \textbf{Formato:} lista de 5 patrones repetidos en tu manera de escribir, cada uno con un ejemplo real tuyo (frase textual). \textbf{Extensión:} 5 patrones + 5 ejemplos. \textbf{Sabes que terminaste cuando} los ejemplos son tuyos (no inventados), específicos (no ``hablo informal''), y un compañero que te conoce podría reconocerte por esos patrones.
\end{infoband}

\puente{Ya identificaste tu voz. Ahora vamos a ver \textbf{qué pasa cuando un texto pasa por la IA} --- qué se conserva, qué se gana, qué se pierde. Es el experimento clave para decidir cómo vas a colaborar con ella.}

\begin{titlebox}{milcTurquesa}
Fase 2 · Sistematización con pensamiento computacional
\end{titlebox}

\begin{softbox}{milcTurquesa}{milcGris}{Cuatro pilares aplicados al tema}
Toda colaboración con IA tiene cuatro pilares. Vamos a descomponerlos para que sepas exactamente cómo dirigir a la IA en lugar de que ella te dirija a ti.
\end{softbox}

\small
\begin{tabularx}{\textwidth}{>{\bfseries\RaggedRight\arraybackslash}p{3.6cm}Y}
\rowcolor{milcTurquesa!90}\color{white}Pilar & \color{white}\bfseries Aplicación al tema\\
1. Descomposición & \textbf{Descomposición.} Una buena colaboración con IA se descompone en 4 capas que actúan como una conversación: (1) \textbf{contexto} (le dices quién eres, qué haces, para qué es el texto, dónde se publica); (2) \textbf{voz} (le compartes 2-3 ejemplos reales de tu manera de escribir --- mensajes, posts antiguos, fragmentos donde te reconoces); (3) \textbf{tarea} (le pides algo concreto y específico, no ``ayúdame''); (4) \textbf{iteración} (revisas la respuesta, marcas lo que funciona, le pides cambios precisos, vuelves a revisar). Sin las 4 capas, sale texto genérico.\\
2. Reconocimiento de patrones & \textbf{Reconocimiento de patrones.} Toda IA generativa sigue el patrón básico ``\emph{le das señal, ella amplía}''. Si tu señal es genérica (``escribe un Sobre mí''), su respuesta será un promedio estadístico de millones de ``Sobre mí'' que vio --- mediocre por construcción. Si tu señal es rica (contexto + ejemplos + tono específico + restricciones explícitas), su respuesta puede acercarse a tu voz. La diferencia entre un texto IA mediocre y uno bueno no está en la IA: está en la calidad del prompt.\\
3. Abstracción & \textbf{Abstracción.} Lo esencial de tu colaboración con IA cabe en una pregunta: \emph{¿estoy usándola como espejo o como autora?}. Si la IA es \textbf{espejo}, tú decides --- ella refleja, amplifica, propone variaciones y tú firmas con criterio. Si la IA es \textbf{autora}, tú firmas algo que no es tuyo, y la diferencia entre las dos posiciones es ética, no técnica. Detectar en qué posición estás cada vez que abres ChatGPT es una habilidad crítica del siglo XXI.\\
4. Algoritmo conceptual & \textbf{Algoritmo.} Estructura mínima de un prompt para coescritura: (1) ``Eres un editor que me ayuda a escribir [tipo de texto]''; (2) ``Aquí están 3 ejemplos reales de mi manera de escribir: [ejemplos]''; (3) ``Necesito que me ayudes a redactar [tarea concreta] de [longitud específica] para [audiencia]''; (4) ``Conserva mi tono. Evita estas palabras: [lista negra]. Dame 3 versiones para elegir.''. Itera con peticiones específicas: ``más corto'', ``menos cliché en el párrafo 2'', ``cambia la metáfora''.\\
\end{tabularx}
\normalsize

\begin{drawbox}{milcTurquesa}{Anatomía de un prompt para coescritura con IA}
\textbf{Rol:} ``Eres un editor de textos personales...'' \\[1mm] \textbf{Voz:} ``Aquí están 3 fragmentos que escribí. Conserva este tono.'' \\[1mm] \textbf{Tarea:} ``Ayúdame a escribir mi `Sobre mí' en 3 párrafos cortos.'' \\[1mm] \textbf{Restricciones:} ``Evita clichés como `apasionado', `creativo', `innovador'.'' \\[1mm] \textbf{Iteración:} ``Dame 3 versiones para elegir una y luego refinarla.''
\end{drawbox}

\begin{softbox}{milcMostaza}{milcGris}{Errores comunes y su corrección}
(1) Pedirle a la IA que escriba sin darle ejemplos de tu voz: sale genérico, suena a manual de ventas estándar. (2) Aceptar la primera versión sin iterar: te quedas con el promedio estadístico, nunca con tu voz. (3) Usar las palabras que la IA mete por defecto (``apasionado'', ``transformador'', ``viaje'', ``aventura'', ``empoderar''): se notan a kilómetros y delatan que el texto es de IA, no tuyo. (4) Firmar como tuyo un texto que no revisaste línea por línea: estás firmando algo que no puedes defender si te preguntan. (5) Pedirle a la IA que ``sea creativa'': la creatividad no se le pide, se le restringe --- ``no uses metáforas obvias'' produce mejor resultado que ``sé creativo''.
\end{softbox}

\begin{infoband}{milcTurquesa}
\textbf{Lo que va al cuaderno (Actividad 2 · ANALIZA).} Título: «Actividad 2 --- Mi voz vs. la IA». \textbf{Formato:} (1) un prompt sencillo (``escribe un Sobre mí para una persona como yo''), pega la respuesta de la IA. (2) Anota 3 cosas que se parecen a tu voz y 3 que no. (3) Reescribe la respuesta a mano corrigiendo lo que no suena a ti. \textbf{Sabes que terminaste cuando} puedes señalar palabras específicas que la IA usó y que tú no usarías nunca.
\end{infoband}

\puente{Tienes el diagnóstico. Toca escribir el texto real: la página ``Sobre mí'' del sitio, en una conversación con la IA donde TÚ diriges y ella ejecuta.}

\begin{titlebox}{milcMagenta}
Fase 3 · Praxis con pensamiento computacional
\end{titlebox}

\begin{softbox}{milcMagenta}{milcGris}{Construyendo el producto con los cuatro pilares}
\textbf{Actividad 3 · CREA --- Tu ``Sobre mí'' coescrito con IA} (35 min · individual). Hora de escribir el texto que va al sitio. Vas a tener una conversación con la IA donde TÚ diriges: le das contexto, ejemplos, tarea, iteras, y al final firmas un texto que puedes defender línea por línea.
\end{softbox}

\small
\begin{tabularx}{\textwidth}{>{\bfseries\RaggedRight\arraybackslash}p{3.6cm}Y}
\rowcolor{milcMagenta!90}\color{white}Pilar & \color{white}\bfseries Aplicación al producto\\
Descomposición & \textbf{Descomposición.} Tu sesión con la IA se descompone en 4 turnos secuenciales: (1) \textbf{prompt inicial} con rol + voz + tarea + restricciones; (2) \textbf{respuesta IA} que lees críticamente subrayando lo que sí suena a ti; (3) \textbf{tu revisión + petición de cambios} específicos (no ``mejóralo'', sí ``el párrafo 2 suena cliché, reemplaza'la metáfora del viaje''); (4) \textbf{versión final que tú reescribes} agregando lo único que la IA no puede saber (una anécdota local, una manía propia).\\
Patrones & \textbf{Patrones.} Sigue el patrón \textbf{``espejo''}: la IA produce, tú diriges. Si en algún momento sientes que estás copiando y pegando sin pensar, detén el flujo --- ya estás dejando que ella sea la autora. La señal de alerta: si lees el texto y no sientes nada particular (ni reconocimiento ni rechazo), es texto IA puro disfrazado de tuyo. Bórralo y vuelve a empezar con más contexto.\\
Abstracción & \textbf{Abstracción.} Cada párrafo expresa una sola idea sobre ti. Si un párrafo de la IA dice 3 cosas distintas (``soy creativo, ordenado y comunicativo''), pártelo o bórralo. Si dice nada concreto (puro adjetivo sin sustento), bórralo. La prueba: cada párrafo debería poder defenderse con una evidencia real --- una pieza, un proyecto, una situación verificable. Sin evidencia, es inflación.\\
Algoritmo de armado & \textbf{Algoritmo.} (1) Abre ChatGPT, Claude o Gemini; (2) prompt inicial completo (rol + voz con ejemplos + tarea + restricciones) usando la plantilla de la Sist; (3) lee la respuesta con marcador en mano y subraya lo que sí suena a ti; (4) pide una segunda versión con esas direcciones específicas; (5) reescribe \emph{a mano} la versión final en el cuaderno --- agregar 1 frase que la IA nunca podría escribir (una referencia a tu barrio, una manía tuya, una anécdota mínima). Esa frase es la firma humana.\\
\end{tabularx}
\normalsize

\begin{drawbox}{milcMagenta}{Checklist antes de cerrar el cuaderno}
\checkbox El texto final cabe en 3 párrafos cortos. \\[1mm] \checkbox Cero palabras de la lista negra (``apasionado'', ``creativo'', ``innovador'', ``transformador''). \\[1mm] \checkbox Hay al menos 1 frase que solo tú podrías escribir (referencia concreta a tu lugar, tu historia, tu manía). \\[1mm] \checkbox Lo lees en voz alta y suena a ti --- ni inflado ni copiado. \\[1mm] \checkbox Coherente con tu propuesta de valor (Guía 1) y sistema visual (Guía 2).
\end{drawbox}

\begin{softbox}{milcMagenta}{milcGris}{Plantilla de trabajo}
\textbf{Párrafo 1 · Quién soy.} \textit{En una frase corta, sin currículum.} \\[1mm] \textbf{Párrafo 2 · Qué hago.} \textit{Tu propuesta de valor en voz propia.} \\[1mm] \textbf{Párrafo 3 · Cómo conectar.} \textit{Email + una invitación cálida y específica.}
\end{softbox}

\begin{infoband}{milcMagenta}
\textbf{Lo que va al cuaderno (Actividad 3 · CREA).} Título: «Actividad 3 --- Sobre mí · versión final». \textbf{Formato:} (1) el prompt completo que usaste con la IA (arriba), (2) el texto final de 3 párrafos cortos (abajo). \textbf{Extensión:} una página de cuaderno. \textbf{Sabes que terminaste cuando} cero palabras de la lista negra, al menos 1 frase irrepetible, y un compañero que te conoce te reconoce al leerlo.
\end{infoband}

\puente{Lo que sigue resume \emph{qué} entregas: una página ``Sobre mí'' lista para pegar en el sitio. El criterio principal: que alguien que te conoce pueda leerla y reconocer tu voz.}

\begin{titlebox}{milcMostaza}
Producto de la sesión
\end{titlebox}
\textbf{Página ``Sobre mí'' coescrita con IA y firmada con tu voz}.
\begin{softbox}{milcMostaza}{milcGris}{Criterios de calidad}
(1) 3 párrafos cortos, no más. (2) Cero palabras clichés (``apasionado'', ``creativo'', ``innovador'', ``transformador''). (3) Al menos 1 frase que solo tú podrías escribir (referencia local, anécdota, manía). (4) Coherente con tu propuesta de valor y sistema visual de las guías anteriores. (5) Pasa el ``read aloud test'': al leerlo en voz alta a alguien que te conoce, te reconoce.
\end{softbox}

\puente{Antes de cerrar, mira el proceso desde las cinco dimensiones humanas. Coescribir con IA es un acto profundamente personal, ciudadano y ético --- aunque parezca solo técnico.}

\begin{titlebox}{milcVino}
Fase 4 · Evaluación liberadora — cinco dimensiones
\end{titlebox}

\begin{softbox}{milcVino}{milcGris}{Sentido de la evaluación}
Evaluar no es solo revisar una nota. Es reconocer qué aprendí, qué cambió en mí, qué emociones manejé, cómo me relaciono con la ciudad y la comunidad, y qué le dejaré a quien viene detrás.
\end{softbox}

\textbf{Mi reflexión liberadora en cinco dimensiones.} Completa cada frase iniciada:

\small
\begin{tabularx}{\textwidth}{>{\bfseries\RaggedRight\arraybackslash}p{3.4cm}Y>{\RaggedRight\arraybackslash}p{4.6cm}}
\rowcolor{milcVino}\color{white}Dimensión & \color{white}\bfseries Mi respuesta (frame starter) & \color{white}\bfseries Algo que descubrí\\
1. Personal & Lo que más cambió en mí fue \linead{6.5cm} & \linead{4.2cm}\\
2. Emocional & La emoción más fuerte fue \linead{2.5cm} y la regulé \linead{3cm} & \linead{4.2cm}\\
3. Ciudadana & En democracia esto importa porque \linead{6cm} & \linead{4.2cm}\\
4. Local (Cartago) & En mi barrio sería útil para \linead{6.5cm} & \linead{4.2cm}\\
5. Intergeneracional & A mi abuela/abuelo le diría \linead{6.5cm} & \linead{4.2cm}\\
\end{tabularx}
\normalsize

\begin{infoband}{milcVino}
\textbf{Para la próxima sustentación.} Vuelve a esta tabla en seis meses. Verás que las respuestas cambian — no porque lo aprendido cambie, sino porque tú cambias. La evaluación liberadora es una conversación contigo a través del tiempo.
\end{infoband}


\puente{Y para terminar, tres voces sobre la autoría con IA. Dussel sobre la voz colonizada, Marco Aurelio sobre la coherencia interna, Floridi sobre la IA como otro inforg en la infoesfera.}

\begin{titlebox}{milcGradeDark}
Triángulo de pensamiento · cierre filosófico
\end{titlebox}

\begin{softbox}{milcVino}{milcGris}{Dussel · lente del nosotros}
\textit{``Toda comunidad de comunicación se hace plena cuando incluye la voz del excluido.''}

La IA fue entrenada principalmente con textos del norte global, en inglés. Su voz por defecto está sesgada hacia ese mundo. Si la dejas escribir sin dirigirla, te dará frases neutras que diluyen tu identidad latinoamericana. Coescribir con criterio es un acto de inclusión de tu propia voz en la infoesfera dominante.

\textbf{¿La IA está incluyendo mi voz local o me la está borrando suavemente? ¿Estoy permitiendo que escriba por mí en lugar de conmigo?}
\end{softbox}
\linead{\linewidth}\\[1mm]
\linead{\linewidth}

\begin{softbox}{milcMostaza}{milcGris}{Estoicismo · Marco Aurelio · cuidado interior}
\textit{``No es lo que dices, es lo que haces lo que define quién eres.''}

La tentación con IA es publicar lo que escribió y firmar como tuyo. Pero firmar es un acto: estás diciéndole al mundo ``esto soy yo''. La integridad estoica te pide que firmes solo lo que puedes defender línea por línea --- y eso depende solo de ti.

\textbf{¿Puedo defender cada línea de este texto como mía? Si me preguntan ``¿por qué pusiste esta frase?'', ¿tengo respuesta?}
\end{softbox}
\linead{\linewidth}\\[1mm]
\linead{\linewidth}

\begin{softbox}{milcTurquesa}{milcGris}{Floridi · lente de la infoesfera}
\textit{``Cuando un agente delega su voz a otro, ambos quedan transformados en la infoesfera.''}

Cada vez que aceptas un texto de la IA sin revisarlo, le delegas tu voz. La IA se convierte en un coautor invisible de tu identidad digital. Floridi llama a eso \emph{re-ontologisation informacional}: cambias quién eres digitalmente al cambiar quién escribe por ti.

\textbf{¿Estoy usando la IA como herramienta o le estoy entregando mi voz? ¿Quién está firmando cuando publico este texto?}
\end{softbox}
\linead{\linewidth}\\[1mm]
\linead{\linewidth}

\begin{titlebox}{milcVino}
Compromiso final
\end{titlebox}

\small
\begin{tabularx}{\textwidth}{>{\RaggedRight\arraybackslash}p{3.0cm}YYY}
\rowcolor{milcVino}\color{white}\bfseries Criterio & \color{white}\bfseries Lo logré & \color{white}\bfseries En proceso & \color{white}\bfseries Necesito apoyo\\
Comprensión del tema & Explico ideas centrales con mis palabras. & Reconozco ideas pero falta precisión. & Necesito repasar conceptos base.\\
Pensamiento computacional & Apliqué los 4 pilares al diseñar mi producto. & Apliqué algunos pilares. & Necesito releer la fase 2.\\
Aplicación y producto & Mi producto cumple los criterios y se puede revisar. & Producto funcional con ajustes pendientes. & Aún en construcción.\\
Reflexión 5 dimensiones & Respondí cada dimensión con honestidad. & Respondí algunas con detalle. & Mis respuestas son muy generales.\\
Triángulo de pensamiento & Conecté las 3 voces con mi trabajo real. & Conecté al menos una. & No relaciono las voces con mi proceso.\\
\end{tabularx}
\normalsize

\textbf{Hoy entendí que:}\\[1mm]
\linead{\linewidth}\\[1mm]
\linead{\linewidth}

\textbf{Mi compromiso para la próxima sesión es:}\\[1mm]
\linead{\linewidth}\\[1mm]
\linead{\linewidth}

\begin{softbox}{milcMostaza}{milcGris}{Cita de despedida · Marco Aurelio}
\textit{``El obstáculo en el camino se vuelve el camino. Lo que estorba al camino se vuelve camino.''}\\[1mm]
Cada error de hoy, bien observado, se convierte en pieza del camino que pisarás mañana.
\end{softbox}

\small
\begin{tabularx}{\textwidth}{>{\bfseries}p{3.6cm}Y}
\rowcolor{milcGradeDark!12}Nombre completo: & \linead{10cm}\\
Fecha: & \linead{4cm} \quad Curso/grupo: \linead{4cm}\\
Mi firma: & \linead{9cm}\\
\end{tabularx}
\normalsize

\begin{infoband}{milcGradeDark}
\textbf{Para la próxima semana.} Antes de publicar cualquier texto coescrito con IA esta semana (redes, correos, ensayos), aplícale el ``read aloud test'': léelo en voz alta a alguien que te conoce. Si te reconoce, publícalo. Si no, reescríbelo. El próximo lunes traes 3 textos pasados por ese filtro.
\end{infoband}

\end{document}
